\section{Latihan}
\label{sec:latihan-bab8}

\begin{enumerate}
    \item \textbf{MOV.} Jelaskan aturan operand pada instruksi
          \asm{MOV}. Sebutkan tiga bentuk pemindahan yang
          diperbolehkan dan dua bentuk yang tidak diperbolehkan.

    \item \textbf{Perbaikan program.} Instruksi
          \lstasm|MOV [BX], [SI]| tidak dapat dirakit. Tuliskan
          penggantinya dengan register perantara, lalu jelaskan
          mengapa cara itu diperlukan.

    \item \textbf{LEA.} Jelaskan perbedaan
          \lstasm|LEA SI, nilai| dan \lstasm|MOV SI, nilai|.
          Kapan \asm{LEA} lebih tepat dipakai?

    \item \textbf{Far pointer.} Diketahui data
          \code{ptr DD 12345678h}.
          \begin{enumerate}
              \item Tuliskan empat byte yang tersimpan di memori.
              \item Tentukan isi \reg{SI} dan \reg{DS} setelah
                    \lstasm|LDS SI, ptr|.
              \item Tentukan isi \reg{DI} dan \reg{ES} setelah
                    \lstasm|LES DI, ptr|.
              \item Jelaskan mengapa \asm{LDS} harus dipakai dengan
                    hati-hati.
          \end{enumerate}

    \item \textbf{XCHG dan CMP.} Diketahui \reg{AX} = \hex{1234}
          dan \reg{BX} = \hex{5678}.
          \begin{enumerate}
              \item Tentukan isi \reg{AX} dan \reg{BX} setelah
                    \lstasm|XCHG AX, BX|.
              \item Tentukan \eng{flag} yang berubah setelah
                    \lstasm|CMP AX, BX|.
              \item Jelaskan apakah \asm{CMP} mengubah isi kedua
                    operand.
          \end{enumerate}

    \item \textbf{XLAT.} Diketahui tabel
          \code{DB '0','1','2','3','4','5','6','7','8','9'},
          \reg{BX} menunjuk awal tabel, dan \reg{AL} = 7.
          \begin{enumerate}
              \item Tentukan isi \reg{AL} setelah \asm{XLATB}.
              \item Jelaskan peran \reg{BX} dan \reg{AL} pada
                    instruksi tersebut.
              \item Sebutkan satu kegunaan \asm{XLAT} dalam program.
          \end{enumerate}

    \item \textbf{CBW dan CWD.} Tentukan hasil setiap instruksi.
          \begin{enumerate}
              \item \reg{AL} = \hex{7F}, lalu \asm{CBW}
              \item \reg{AL} = \hex{80}, lalu \asm{CBW}
              \item \reg{AX} = \hex{7FFF}, lalu \asm{CWD}
              \item \reg{AX} = \hex{8000}, lalu \asm{CWD}
          \end{enumerate}
          Jelaskan mengapa hasilnya berbeda untuk bilangan positif
          dan negatif.

    \item \textbf{Penulisan program.} Tulislah program yang:
          \begin{enumerate}
              \item mendefinisikan dua variabel word;
              \item menyalin keduanya ke register;
              \item menukar isi kedua register dengan \asm{XCHG};
              \item membandingkan keduanya dengan \asm{CMP};
              \item menyimpan hasil pertukaran ke variabel.
          \end{enumerate}
\end{enumerate}

\begin{refleksibox}
\vspace{1.1em}
Jawab pertanyaan berikut pada jurnal praktikum Anda.
\begin{enumerate}
    \item Instruksi mana yang paling sulit dipahami: \asm{LEA},
          \asm{LDS/LES}, \asm{XLAT}, atau \asm{CBW/CWD}? Mengapa?
    \item Setelah menelusuri program, apakah Anda dapat membedakan
          dengan jelas antara alamat dan isi memori? Jelaskan.
    \item Konsep apa yang perlu Anda ulangi sebelum mempelajari
          interupsi BIOS/DOS pada bab berikutnya?
\end{enumerate}
\end{refleksibox}
